%\addtolength{\hoffset}{-2.25cm}
%\addtolength{\textwidth}{4.5cm}
%\addtolength{\voffset}{-3.25cm}
%\addtolength{\textheight}{10cm}

%\setlength{\parskip}{5pt}
%\setlength{\parindent}{0in}

\usepackage[left=2.54cm, right=2.54cm,
bottom=2.54cm,
  top=2cm, % Establece el espacio superior
  headheight=2cm, % Establece la altura del encabezado
  headsep=0.5cm, % Establece el espacio entre el encabezado y el cuerpo del texto
]{geometry}



%----------------------------------------------------------------------------------------
%	PACKAGES AND OTHER DOCUMENT CONFIGURATIONS
%----------------------------------------------------------------------------------------
\usepackage{blindtext} % Package to generate dummy text
%\usepackage{lmodern} % Use the Charter font
%\usepackage[utf8]{inputenc} % Use UTF-8 encoding
\usepackage{microtype} % Slightly tweak font spacing for aesthetics
\usepackage[spanish,mexico,es-lcroman]{babel} % Language hyphenation and typographical rules
\usepackage{multirow}
\usepackage[table,xcdraw]{xcolor}
\usepackage{caption}
\usepackage{enumerate} % enumerado

%\usepackage{hyperref}
\spanishdecimal{.}
\usepackage[most]{tcolorbox}
\tcbuselibrary{listingsutf8}

\usepackage{tikz,tcolorbox}
\usepackage[framemethod=tikz]{mdframed}

\usepackage{framed}
\usepackage{color}
\usepackage{wrapfig}
\definecolor{shadecolor}{RGB}{0.86,0.84,0.82}

\usepackage{amsthm, amsmath, amssymb} % Mathematical typesetting
\usepackage{amsfonts}
\usepackage{float} % Improved interface for floating objects
\usepackage[final, colorlinks = true, 
            linkcolor = black,
            %linktoc = false,
            citecolor = blue,
            urlcolor = blue!80!black!70]{hyperref} % For hyperlinks in the PDF
%\usepackage{hyperref}
%\hypersetup{
  %colorlinks=true,
  %linkcolor=black,
  %urlcolor=blue
%}
\usepackage{titletoc}
\usepackage{titlesec}

\usepackage{graphicx, multicol} % Enhanced support for graphics
\usepackage{xcolor} % Driver-independent color extensions
\usepackage{marvosym, wasysym} % More symbols
\usepackage{rotating} % Rotation tools
\usepackage{censor} % Facilities for controlling restricted text
\usepackage{listings} % Environment for non-formatted code, !uses style file!
\usepackage{pseudocode} % Environment for specifying algorithms in a natural way
\usepackage{booktabs} % Enhances quality of tables
\usepackage{tikz-qtree} % Easy tree drawing tool
\tikzset{every tree node/.style={align=center,anchor=north},
         level distance=2cm} % Configuration for q-trees
%----------------- ACA
%\usepackage[backend=biber,style=numeric,
            %sorting=nyt]{biblatex} % Complete reimplementation of bibliographic facilities
%\addbibresource{ecl.bib}


\usepackage[autostyle]{csquotes} % Context sensitive quotation facilities
\usepackage[yyyymmdd]{datetime} % Uses YEAR-MONTH-DAY format for dates
\renewcommand{\dateseparator}{-} % Sets dateseparator to '-'

%LETRA MODIFICADA ---------------
\usepackage{lmodern}
\renewcommand*\familydefault{\sfdefault} % use lmodern for paragraphs
\usepackage[T1]{fontenc} % i.a. makes title font non-serif

\usepackage{fouriernc} % sharper font for on-screen
\DeclareMathAlphabet{\mathbf}{T1}{fut\mathfamilyextension}{b}{n} 
\renewcommand{\bfdefault}{b}

\newcommand{\note}[1]{\marginpar{\scriptsize \textcolor{red}{#1}}} % Enables comments in red on margin
\usepackage{array}

\newcommand*{\filltext}{The quick brown fox jumps over the lazy dog.}

%\usepackage{adjustbox}
\usepackage{graphicx,bm}
%\usepackage{showframe}   %% just for demo

\usepackage{booktabs}% http://ctan.org/pkg/booktabs
\newcommand{\tabitem}{~~\llap{\textbullet}~~}
\usepackage{mathtools}
\usepackage{array}
\newcolumntype{L}{>{\centering\arraybackslash}m{7cm}}

\usepackage[font={small,it},labelfont={bf}]{caption}
\newcommand{\centered}[1]{\begin{tabular}{l} #1 \end{tabular}}

\definecolor{bl1}{RGB}{170, 170, 170}
\definecolor{color1}{RGB}{60, 93, 189}
\definecolor{color2}{RGB}{56, 107, 255}
\definecolor{color3}{RGB}{0, 48, 188}
\usepackage{nicematrix}

\usepackage{tikz}
\newcommand*\circled[1]{\tikz[baseline=(char.base)]{\node[shape=circle,draw,inner sep=2pt] (char) {#1};}}

\usepackage{setspace}

%\renewcommand{\familydefault}{\sfdefault}
%\usepackage{helvet}

\usepackage{tabularray}
\usepackage[export]{adjustbox}

\UseTblrLibrary{varwidth}
\usepackage{soul}
%\usepackage[%
  %colorlinks=true,%
  %urlcolor=red%
%]{hyperref}
\usepackage[pages=some]{background}

\usepackage{tabularray,varwidth,enumitem}

%\usepackage{titlesec}
%\newcommand{\chapfnt}{\fontsize{12}{12}}
\newcommand{\secfnt}{\fontsize{12}{12}}
\newcommand{\ssecfnt}{\fontsize{12}{12}}

%PAQUETE PARA AGREGAR + A LAS FLECHAS
\usepackage{amssymb,stackengine,scalerel}

%SIMBOLO +/- REDUCIDO
\newcommand{\MAS}{\mbox{\scriptsize+}}
\newcommand{\MENOS}{\mbox{\footnotesize-}}
%SENTIDOS POSITIVOS
%SENTIDO ANTIHORARIO CON SUMA
\newcommand\plusarrowleft{\mathrel{\ThisStyle{%
  \ensurestackMath{\stackengine{-0.5pt}{\SavedStyle\curvearrowleft}%
        {\SavedStyle\hspace{2.5pt}\MAS}{O}{c}{F}{T}{L}}%
}}}
%SENTIDO DERECHA CON SUMA
\newcommand\plusrightarrow{\mathrel{\ThisStyle{%
  \ensurestackMath{\stackengine{4.5pt}{\SavedStyle\rightarrow}%
        {\SavedStyle\hspace{-2pt}\MAS}{O}{c}{F}{T}{L}}%
}}}
%SENTIDO ARRIBA CON SUMA
\newcommand\plusuparrow{\mathrel{\ThisStyle{%
  \ensurestackMath{\stackengine{0pt}{\SavedStyle\uparrow}%
        {\SavedStyle\hspace{-7pt}\MAS}{O}{c}{F}{T}{L}}%
}}}
%SENTIDOS NEGATIVOS
%SENTIDO HORARIO CON RESTA
\newcommand\minarrowright{\mathrel{\ThisStyle{%
  \ensurestackMath{\stackengine{-0.5pt}{\SavedStyle\curvearrowright}%
        {\SavedStyle\hspace{-1.5pt}\MENOS}{O}{c}{F}{T}{L}}%
}}}
%SENTIDO IZQUIERDA CON RESTA
\newcommand\minleftarrow{\mathrel{\ThisStyle{%
  \ensurestackMath{\stackengine{4.5pt}{\SavedStyle\leftarrow}%
        {\SavedStyle\hspace{2pt}\MENOS}{O}{c}{F}{T}{L}}%
}}}
%SENTIDO ABAJO CON RESTA
\newcommand\mindownarrow{\mathrel{\ThisStyle{%
  \ensurestackMath{\stackengine{2pt}{\SavedStyle\downarrow}%
        {\SavedStyle\hspace{-7pt}\MENOS}{O}{c}{F}{T}{L}}%
}}}

\usepackage{amssymb}% http://ctan.org/pkg/amssymb
\usepackage{pifont}% http://ctan.org/pkg/pifont
\newcommand{\cmark}{\ding{51}}%
\newcommand{\xmark}{\ding{55}}%


%CONFIGURACION PARA LONGTABLE
\DefTblrTemplate{caption}{default}{%
  \UseTblrTemplate{caption-tag}{default}%
  \UseTblrTemplate{caption-sep}{default}%
  \UseTblrTemplate{caption-text}{default}%
}
\DefTblrTemplate{capcont}{default}{%
  \UseTblrTemplate{caption-tag}{default}%
  \UseTblrTemplate{caption-sep}{default}%
  \UseTblrTemplate{caption-text}{default}
  \UseTblrTemplate{conthead-text}{default}%
}

\SetTblrStyle{caption-text}{font=\\[-10pt]\it\footnotesize}
\SetTblrStyle{caption-tag}{font=\footnotesize}
%\SetTblrStyle{caption-tag}{font=\bfseries\footnotesize}

\DefTblrTemplate{contfoot-text}{normal}{\scriptsize\textit{Continua en la siguiente página}}
\SetTblrTemplate{contfoot-text}{normal}

\DefTblrTemplate{conthead-text}{default}{\footnotesize(continuación)}

\DefTblrTemplate { caption-sep } { empty } { }
\DefTblrTemplate { caption-sep } { normal } { \vspace*{6pt} }
\SetTblrTemplate { caption-sep } { normal }

\usepackage{nicefrac,xfrac}


\newcommand{\mycbox}[1]{\tikz{\path[draw=black,fill=#1] (0,0) rectangle (0.35cm,0.18cm);}}

% :::::::::::::::: SUBFIGURAS :::::::::::::::::
\usepackage{subcaption}
\captionsetup{subrefformat=parens}
\usepackage{cancel}
% ------------
%\usepackage{newtxtext}
\usepackage[labelsep=period]{caption}
\captionsetup[figure]{font=footnotesize,labelfont={it}}
\captionsetup[sub]{font=normalsize, labelfont=normal}

%----------------------------------------------------------------------------------------

%NUEVO PARRAFO SIN NUMERACION
\titleformat*{\paragraph}{\itshape\bfseries}
\newcommand{\parrafo}[1]{\paragraph{#1}\mbox{}\vspace{6pt}\\}

\setcounter{secnumdepth}{3}
%SÍMBOLO DE ESTRIBO 
\newcommand{\ESTRIBO}{%
  \begingroup\normalsize
  \includegraphics[height=\fontcharht\font`\B]{img/Estribo.pdf}%
  \endgroup
}
\newcommand{\ESTRIBOX}{%
  \begingroup\scriptsize
  \includegraphics[height=\fontcharht\font`\B]{img/Estribo.pdf}%
  \endgroup
}

%SIMBOLO DE ARROBA
\newcommand{\ARR}{\fontfamily{ptm}\selectfont{@}}
%\usepackage{txfonts}


%PARA TEXTO SUBESCRITO
\usepackage{fixltx2e}
\definecolor{verde0}{RGB}{0,170,73}

%ENTORNO PARA CUADRO DE TEXTO

\definecolor{fondo}{RGB}{200,200,200}
\definecolor{azul1}{RGB}{36,22,168}
\definecolor{azul2}{RGB}{9,0,168}
\definecolor{azul3}{RGB}{0,46,211}

% NUEVO ESTILO PARA LA TABLA DE CONTENIDOS

\usepackage{tocloft}

\renewcommand{\cftsecleader}{\bfseries\color{black}\cftdotfill{\cftdotsep}}
\renewcommand{\cftsecaftersnum}{\hspace{1em}\vrule height 1.5ex depth 0pt width 0.1em\hspace{1em}}
\setlength{\cftsecnumwidth}{3em}
\renewcommand{\cftsecfont}{\normalsize\bfseries\hypersetup{linkcolor=black}}
\renewcommand{\cftsecpagefont}{\color{black}\normalsize\bfseries}

%::::::: PARA EL ÍNDICE GENERAL :::::::::::
\renewcommand{\cfttoctitlefont}{\color{black}\Large\bfseries\MakeUppercase}
\renewcommand{\cftaftertoctitle}{}
\renewcommand{\cftaftertoctitleskip}{0pt}

% Añadir líneas horizontales alrededor del título
\renewcommand{\cftaftertoctitle}{%
  \vspace{6pt}
  \color{black}\hrule height 1pt\relax
  \vspace{6pt}%
  \color{black}%
}

%::::::: PARA EL ÍNDICE DE FIGURAS :::::::::::
\renewcommand{\cftloftitlefont}{\color{black}\Large\bfseries\MakeUppercase}
\renewcommand{\cftafterloftitle}{}
\renewcommand{\cftafterloftitleskip}{10pt}
%
\renewcommand{\cftafterloftitle}{%
  \vspace{6pt}
  \color{black}\hrule height 1pt\relax
  \vspace{6pt}%
  \color{black}%
}


%::::::: PARA EL ÍNDICE DE TABLAS :::::::::::
\renewcommand{\cftlottitlefont}{\color{black}\Large\bfseries\MakeUppercase}
\renewcommand{\cftafterlottitle}{}
\renewcommand{\cftafterlottitleskip}{10pt}
%
\renewcommand{\cftafterlottitle}{%
  \vspace{6pt}
  \color{black}\hrule height 1pt\relax
  \vspace{6pt}%
  \color{black}%
}



%\usepackage[natbib=true,backend=biber,sorting=nyt,style=apa6]{biblatex}

%\renewcommand*{\bibfont}{\fontsize{12}{12}\selectfont}
%\addbibresource{Bibliografia.bib}

\renewcommand{\tablename}{Tabla}
\definecolor{ROJO1}{RGB}{255, 18, 58}
\definecolor{AMARILLO1}{RGB}{231, 255, 54}
\definecolor{VERDE1}{RGB}{0, 255, 91}
\definecolor{VERDE2}{RGB}{0, 132, 53}

%\definecolor{verde3}

\usepackage{fancyhdr} % Headers and footers
%Nuevo comando para definir la seccion

%\newcommand{\miseccion}{}
%\renewcommand{\sectionmark}[1]{\renewcommand{\miseccion}{\thesection.\enspace#1}}

%\newif\ifappendix

% Define un nuevo comando \lastsectionmark para almacenar el nombre de la última sección

\newcommand{\miseccion}{}
\renewcommand{\sectionmark}[1]{\renewcommand{\miseccion}{\thesection.\enspace#1}}

%\newcommand{\miseccion}{}
%\renewcommand{\sectionmark}[1]{%\renewcommand{\miseccion}{\rightmark}%}


\fancypagestyle{estilo1}{
  \fancyhf{}
  \renewcommand{\headrulewidth}{0pt} % Elimina la línea horizontal en la parte superior
  \fancyhead[R]{\roman{\thepage}} % Coloca el número de página en romanos minúsculos en la esquina superior derecha la derecha
}

\fancypagestyle{estilo2}{
  \fancyhf{}
  \renewcommand{\headrulewidth}{0pt}
  \fancyhead[L]{}
  \fancyhead[R]{} % La numeración de página a la derecha
  \setcounter{page}{1}
  \fancyfoot[L]{}
  \fancyfoot[C]{-- \thepage\ --}
  \fancyfoot[R]{}
}


\usepackage[toc,page]{appendix}
\renewcommand\appendixname{Anexos}
\renewcommand\appendixtocname{\appendixname}
\renewcommand\appendixpagename{Anexos}


%\usepackage[natbib=true,backend=biber,sorting=nyt,style=ieee]{biblatex}
\usepackage[backend=biber,style=ieee,sorting=nyt]{biblatex}
\addbibresource{Referencias.bib}

\makeatletter
\DeclareCiteCommand{\supercite}[\mkbibsuperscript]
  {\usebibmacro{cite:init}}
  {\usebibmacro{citeindex}%
   \mkbibbrackets{\usebibmacro{cite:comp}}}
  {}
  {\usebibmacro{cite:dump}}
\makeatother


\usepackage{makecell}
\captionsetup[subtable]{position=auto, justification = centering}

%ETIQUETA DE TABLA
\DeclareCaptionLabelSeparator*{spaced}{\\[1pt]}
\captionsetup[table]{font=footnotesize, textfont=it,format=plain,justification=justified, singlelinecheck=false,labelsep=spaced,skip=6pt, labelfont={normal,footnotesize}}

%NOTAS AL PIE DE LA TABLA
\usepackage[flushleft]{threeparttable}


\tcbset{%
    theo/.style={%
        enhanced,
        breakable,
        sharp corners,
        toprule=0pt, rightrule=0pt, bottomrule=0pt, leftrule=1mm,
        colback=#1!5, colframe=#1!80!black, coltitle=#1!80!black, 
        detach title,
        overlay unbroken and first ={
            \node[rotate=90, minimum width=1cm, anchor=south, font=\bfseries] 
               at (frame.west) {\tcbtitle};
        }
    }
}

\newtcbtheorem[no counter]{mytheo}{}
{theo=gray, fonttitle=\scriptsize, separator sign none}{th}


%Enumeracion exclusiva para ANEXOS 
\usepackage{chngcntr}
%\newcommand{\anexosection}[1]{%
  %\refstepcounter{section}%
  %\renewcommand{\thesection}{\Alph{section}}%
  %\section{#1}%
%}

%ESTILO DE FUENTE ROBOTO
\newcommand\TITULOS{\fontfamily{qag}\selectfont}


%ESTILO TECLAS 
\usepackage{inconsolata}
\definecolor{micolor1}{RGB}{244,242,249}
\definecolor{micolor2}{RGB}{000,114,190}
\newtcbox{\TECLAS}[1][red]{on line,
arc=3pt,colback=blue!8,colframe=blue!8,
before upper={\rule[-3pt]{0pt}{10pt}},boxrule=1pt,
boxsep=0pt,left=2pt,right=2pt,top=1pt,bottom=.5pt}

%NUMEROS CIRCULARES
\makeatletter
% Based on https://tex.stackexchange.com/questions/7032/good-way-to-make-textcircled-numbers
%
\usepackage{pict2e,picture}
\pgfkeys{/csteps/inner ysep/.initial=4pt,
    /csteps/inner xsep/.initial=4pt,
    /csteps/inner color/.initial=red,
    /csteps/outer color/.initial=blue,
}
\newsavebox\csteps@CBox
\newlength\csteps@XLength \newlength\csteps@YLength \newlength\csteps@YDepth \newlength\csteps@tmplen
\def\csteps@CircledParam#1#2{\sbox\csteps@CBox{#2}%
    \csteps@XLength=\wd\csteps@CBox\advance\csteps@XLength by\pgfkeysvalueof{/csteps/inner xsep}\relax
    \csteps@tmplen=\pgfkeysvalueof{/csteps/inner ysep}\relax
    \csteps@YDepth=\dp\csteps@CBox\advance\csteps@YDepth by 0.5\csteps@tmplen\relax
    \csteps@YLength=\ht\csteps@CBox\advance\csteps@YLength by\dp\csteps@CBox\advance\csteps@YLength by\pgfkeysvalueof{/csteps/inner ysep}\relax
    \typeout{DBG:#2\space X\space\the\csteps@XLength\space Y:\the\csteps@YLength\space D:\the\csteps@YDepth}%
    \raisebox{-#1\csteps@YDepth}{%
    \ifdim\csteps@XLength>\csteps@YLength
    \makebox[\csteps@XLength]{% X bigger than Y
        \makebox(0,\csteps@YLength){%
            \color{\pgfkeysvalueof{/csteps/outer color}}\put(0,0){\oval(\csteps@XLength,\csteps@YLength)}%
        }%
    \makebox(0,\csteps@YLength){%
        \put(-.5\wd\csteps@CBox,0){\textcolor{\pgfkeysvalueof{/csteps/inner color}}{#2}}%
    }}%
    \else
    \makebox[\csteps@YLength]{%
        \makebox(0,\csteps@YLength){%
            \color{\pgfkeysvalueof{/csteps/outer color}}\put(0,0){\circle{\csteps@YLength}}%
        }%
    \makebox(0,\csteps@YLength){%
        \put(-.5\wd\csteps@CBox,0){\textcolor{\pgfkeysvalueof{/csteps/inner color}}{#2}}%
     }}%
    \fi
    }%
}
\def\Circled#1{\csteps@CircledParam{1}{#1}}
\def\CircledTop#1{\csteps@CircledParam{0}{#1}}
\makeatother

%PAQUETE PARA ALGORITMO
\usepackage{etoolbox}
\usepackage[linesnumbered,ruled,vlined]{algorithm2e}

\makeatletter
%%% Coloring the comment as blue
\newcommand\mycommfont[1]{\tiny\ttfamily\textcolor{blue}{#1}}
\SetCommentSty{mycommfont}
\SetKwInput{KwInput}{Entrada}                % Set the Input
\SetKwInput{KwOutput}{Salida}              % set the Output
\SetKwInput{KwData}{Datos}

%\SetAlFnt{\scriptsize}
%\SetAlCapFnt{\scriptsize}
%\SetAlCapNameFnt{\scriptsize}
\SetAlFnt{\tiny}
\SetAlCapFnt{\tiny}
\SetAlCapNameFnt{\tiny}
%\DeclareCaptionLabelFormat{algnonumber}{Algoritmo}
%\captionsetup[algorithm]{labelformat=algnonumber}
\SetAlgorithmName{Algoritmo}{Mi Lista de algoritmos}
%\SetKwFor{While}{Mientras que}
%\SetKwFor{MientrasQue}{mientras que}{hacer}{fin mientras que}
\makeatother
\renewcommand{\thealgocf}{}


\definecolor{titlepagecolor}{RGB}{248,248,248}
\definecolor{namecolor}{RGB}{130,130,130} 



\tcbuselibrary{skins}
\usetikzlibrary{shadings}
\newcounter{datos}
%ESTILO DE CUADRO
\tcbset{
  base/.style={
    empty,
    frame engine=path,
    colframe=gray!4,
    sharp corners,
    title={},  % Eliminamos el título predeterminado
    attach boxed title to top left={yshift*=-\tcboxedtitleheight},
    boxed title style={size=minimal, top=4pt, left=4pt},
    coltitle=black,
    fonttitle=\bfseries\sffamily,
  }
}

\newtcolorbox[use counter=datos]{DATOS}[2][]{%
  base,
  title={#2},  % Usamos el título proporcionado o ninguno si no se proporciona
  boxed title style={
    overlay={
      \draw[black,line width=3pt] (frame.north west)--(frame.north east);
    }
  },
  colback=black,
  overlay unbroken={
    \draw[black] ([yshift=-1.5pt]title.north east)--([xshift=-0.5pt, yshift=-1.5pt]title.north-|frame.east);
  },
  #1  % Posibilidad de pasar opciones adicionales al cuadro
}


\newcommand\ssbullet[1][.8]{\mathbin{\ThisStyle{\vcenter{\hbox{%
  \scalebox{#1}{$\SavedStyle\bullet$}}}}}%
}
\renewcommand{\labelitemi}{\(\ssbullet\)}

\titleformat{\section}
  {\normalfont\Large\bfseries}
  {\thesection.}
  {.5em}
  {}
  [\vspace{.2ex}\titlerule]

\captionsetup[subtable]{position=auto, justification = centering}

%PARA EL ESPECTRO
\usepackage{tkz-fct}
\UseTblrLibrary{diagbox}
\usepackage{scrextend}

\usepackage{pgfplots}
\usepgfplotslibrary{colorbrewer}
\usepgfplotslibrary{colormaps}

\usepackage{pgfgantt}
\usepackage{rotating}

\tikzset{
    pline/.style={% plot line
        every plot/.style={
            mark=x,
            mark options={
                black,
                thin,
                smooth,
            },
            mark size=2pt,
        },
        thin,
        black,
        %dashed,
    },
}
\newcommand{\rectangle}{{%
  \ooalign{$\sqsubset\mkern3mu$\cr$\mkern3mu\sqsupset$\cr}%
}}

% ADORNOS PARA TEXTO
\usepackage{fourier-orns}
\usetikzlibrary{arrows.meta}

%ESTILO HISTERETICO PARA GRAFICOS
\usetikzlibrary{decorations.markings}
%\pgfplotsset{compat=1.13}
\tikzset{
  set arrow inside/.code={\pgfqkeys{/tikz/arrow inside}{#1}},
  set arrow inside={end/.initial=>, opt/.initial=},
  /pgf/decoration/Mark/.style={
    mark/.expanded=at position #1 with
    {
      \noexpand\arrow[\pgfkeysvalueof{/tikz/arrow inside/opt}]{\pgfkeysvalueof{/tikz/arrow inside/end}}
    }
  },
  arrow inside/.style 2 args={
    set arrow inside={#1},
    postaction={
      decorate,decoration={
        markings,Mark/.list={#2}
      }
    }
  },
}

\pgfplotsset{
    din xlabel/.style={
        xticklabel style={
            name=label\ticknum,
            append after command=\pgfextra{\xdef\lastticknum{\ticknum}}
        },
        after end axis/.code={
            \pgfmathparse{int(\lastticknum-1)}
            \path (label\lastticknum.base) -- (label\pgfmathresult.base) node [midway, anchor=base] {#1};
        }
    }
}

\newtcolorbox{SISMO}[1]%
  {
    enhanced,
    sharp corners,
%arc=1mm,
outer arc=1mm,
boxrule=0mm,
toprule=0mm,
bottomrule=0mm,
left=1mm,
right=1mm,
leftrule=0mm,
titlerule=0mm,
toptitle=0mm,
bottomtitle=0mm,
top=0mm,
boxrule = 0.75pt,
colframe=black!55!white,
colback=black!2!white,
title=#1,
fonttitle=\bfseries\sffamily\normalsize,fontupper=\normalsize\itshape,
 coltitle=black,
  attach boxed title to top left={
                                xshift=0mm,
                                yshift=-5.5mm,
                                yshifttext=-1mm
                                },  
top=1.5em,
boxed title style={
                      sharp corners,
                      size=small,
                      colback=gray!10,
                    colframe=black!45!white,colbacktitle=black
                    } 
  }


\begin{comment}
   \titleformat{\section}
  {\normalfont\Large\bfseries}{\thesection \quad {#1}}{.5em}{\vspace{.5ex}}[\titlerule]
\titlespacing*{\section}
  {0pt}{0pt}{15pt} 
\end{comment}
